\documentclass[onecolumn, draftclsnofoot,10pt, compsoc]{IEEEtran}
\usepackage{graphicx}
\usepackage{setspace}
\usepackage{comment}
\usepackage{bigstrut}
\usepackage{geometry}
\usepackage{tabu}
\usepackage{supertabular}
\usepackage{hyperref}
\usepackage{url}
\hypersetup{
  colorlinks=true, linkcolor=blue, citecolor=blue, filecolor=blue, urlcolor=blue}
\geometry{textheight=9.5in, textwidth=7in}

% 1. Fill in these details
\def \CapstoneTeamName{		Remix}
\def \CapstoneTeamNumber{		61}
\def \GroupMemberOne{			Josh Matteson}
\def \GroupMemberTwo{			Steven Powers}
\def \GroupMemberThree{			Evan Tschuy}
\def \CapstoneProjectName{		Many Voices Platform}
\def \CapstoneSponsorCompany{	Oregon State University}
\def \CapstoneSponsorPerson{		Carlos Jensen}

% 2. Uncomment the appropriate line below so that the document type works
\def \DocType{
				%Abstract
		        %Problem Statement
				%Requirements Document
				Technology Review
				%Design Document
				%Progress Report
				}

\newcommand{\NameSigPair}[1]{\par
\makebox[2.75in][r]{#1} \hfil 	\makebox[3.25in]{\makebox[2.25in]{\hrulefill} \hfill		\makebox[.75in]{\hrulefill}}
\par\vspace{-12pt} \textit{\tiny\noindent
\makebox[2.75in]{} \hfil		\makebox[3.25in]{\makebox[2.25in][r]{Signature} \hfill	\makebox[.75in][r]{Date}}}}
% 3. If the document is not to be signed, uncomment the RENEWcommand below
\renewcommand{\NameSigPair}[1]{#1}

%%%%%%%%%%%%%%%%%%%%%%%%%%%%%%%%%%%%%%%
\begin{document}
\begin{titlepage}
    \pagenumbering{gobble}
    \begin{singlespace}
    	\includegraphics[height=4cm]{coe_v_spot1}
        \hfill
        % 4. If you have a logo, use this includegraphics command to put it on the coversheet.
        %\includegraphics[height=4cm]{CompanyLogo}
        \par\vspace{.2in}
        \centering
        \scshape{
            \huge CS Capstone \DocType \par
            {\large\today}\par
            \vspace{.5in}
            \textbf{\Huge\CapstoneProjectName}\par
            \vfill
            {\large Prepared for}\par
            \Huge \CapstoneSponsorCompany\par
            \vspace{5pt}
            {\Large\NameSigPair{\CapstoneSponsorPerson}\par}
            {\large Prepared by }\par
            Group\CapstoneTeamNumber\par
            % 5. comment out the line below this one if you do not wish to name your team
            \CapstoneTeamName\par
            \vspace{5pt}
            {\Large
                \NameSigPair{\GroupMemberOne}\par
                \NameSigPair{\GroupMemberTwo}\par
                \NameSigPair{\GroupMemberThree}\par
            }
            \vspace{20pt}
        }
        \begin{abstract}
        % 6. Fill in your abstract
		\noindent The Many Voices Publishing Platform uses a variety of technologies to handle different
		aspects of the project, from the user interface to the backend database operations. These
		technologies enable to the Many Voices Publishing Platform to succeed in delivering a working
		platform for textbook collaboration.
        \end{abstract}
    \end{singlespace}
\end{titlepage}
\newpage
\pagenumbering{arabic}
\tableofcontents
% 7. uncomment this (if applicable). Consider adding a page break.
%\listoffigures
%\listoftables
\clearpage


% 8. now you write!
\section{Revision Log}
\begin{flushleft}
\tablehead{}
\begin{supertabular}{|p{3cm}|p{3cm}|p{3cm}|p{7cm}|}
\hline
Name & Change Number & Date & Description of Change
\\\hline
Steven Powers & 1 & 2/13/2017 & \begin{itemize}
								\item{Updated document to new design.}
								\item{Changed selection from Angular2 to Aurelia for User Interface Tools.}
								\item{Increased Readability of PDF and Latex document.}
								\item{Drastically reduced complexity of document,
									and slimmed down the document from nearly 36 pages to 22
									pages by setting better page flow practices.}
								\end{itemize}
\\ \hline

Evan Tschuy & 2 & 2/17/2017 & \begin{itemize}
								\item{Changed out HTML option for Latex option.}
								\end{itemize}
\\ \hline

\end{supertabular}
\end{flushleft}


\bigskip

\section{Introduction}


\noindent The Many Voices Publishing Platform is being developed for the purpose of fixing the problems currently associated with the textbook market. We will accomplish this by giving the MVP Platform an easy to use interface, a search bar with a built in results pane, source control, and many other features. Authorship is divided by subsection header.


\newpage

\section{Steven Powers}


{\subsection{User Interface Tools}
\noindent Option 1 - React \cite{React} \\
\noindent React is a JavaScript rendering engine that is developed by Facebook. Originally used with Instagram, React is often paired with Redux for added functionality. React is a popular JavaScript library meant for building user interfaces that is component based.

\medskip

\noindent Option 2 - Aurelia \cite{Aurelia} \\
\noindent Aurelia is a newer JavaScript client framework for mobile, desktop, and the web, by using simplistic integration.

\medskip

\noindent Option 3 - Ember \cite{Ember} \\
\noindent Ember uses web components and templates to increase productivity.

\medskip

\noindent Option 4 - Angular2 \cite{Angular2} \\
\noindent Angular2 is a project started by Google for their internal Green Tea project. Angular2 is a widely documented JavaScript cross-platform library that is used to create native mobile and desktop web applications.

\medskip
\subsubsection{Goals}
{\noindent The use of this technology will aid in the development of the user interface. Having a beautiful and scalable user interface will help users interact with the platform more easily, on whatever device they choose to use it on. \par}

\medskip

\subsubsection{Evaluation Criteria}
\noindent The options are evaluated on

\begin{itemize}
\item Ease of Use
\item File Size
\item Features
\item Performance
\item Standards Compliance
\item Non-Compliance
\item Release
\item License
\end{itemize}



\newpage
\subsubsection{Option Comparison} ~\\

\tablehead{}
\begin{supertabular}{|p{2.5cm}|p{3.3cm}|p{3.3cm}|p{3.3cm}|p{3.3cm}|}
\hline
	\cite{Eisenberg}
	& React
	& Aurelia
	& Ember
	& Angular 2 \\
\hline
	Ease of Use
	& Substantial setup required for working system, lots of documentation and tutorials.
	& Simple setup using NPM and installation
	& Simple setup using NPM and installation
	& Substantial setup required for working system, lots of documentation and tutorials. \\
\hline
	File Size
	& 156kb to ???kb, due to added frameworks
	& 323kb & 435kb
	& 1023kb \\
\hline
	Features
	& View rendering engine with plugin frameworks
	& Router, Animation, HTTP Client & Router, HTTP Client
	& Router, HTTP Client \\
\hline
	Performance (Paints \hspace{3em} per Second)
	& 45-50 & 90-150 (Higher end with additional plugins)
	& 60-100 & 80-130 (Higher end with additional plugins) \\
\hline
	Standards \hspace{3em} Compliance
	& ES 2015
	& HTML, ES 2016, Web Components
	& HTML, ES 2015
	& ES 2016\\
\hline
	Non-Compliance
	& JSX
	& N/A
	& N/A
	& NG2 Markup, Dart \\
\hline
	Release
	& 15.x
	& Beta
	& 2.x
	& Release Candidate \\
\hline
	License
	& BSD
	& MIT
	& MIT
	& MIT \\
\hline
\end{supertabular}

\medskip

\subsubsection{Discussion}
\noindent All of the chosen options have their pros and cons for our web application. All of them however would be a learning and research experience. Angular2 and React have the benefit of being created by large software companies, Google and Facebook respectively. This means that there will be large adoption and documentation / tutorials available. Aurelia and Ember seem to be easier to implement however, they are much newer products and they have a smaller adoption population. This could prove troublesome if we run into problems.

\medskip
\subsubsection{Selection}
\noindent Initially we were set on using Angular2 as part of the team has experience using this JavaScript library, before meeting with our client. Angular2 has a wide adoption and is used by Google for internal projects so the longevity of the framework is expected to last. While the wider range of knowledge backing Angular2 would be helpful, we plan to use Aurelia as our JavaScript framework for our user interface.



\newpage
\subsection{User Login \& Authentication}
\noindent Option 1 - OpenID Connect \\
\noindent OpenID Connect allows for clients of all types, including browser-based JavaScript and native mobile apps, to launch sign-in flows and receive verifiable assertions about the identity of signed-in users \cite{OpenID}.

\medskip

\noindent Option 2 - Facebook  \\
\noindent Facebook Login for Apps is a fast and convenient way for people to create accounts and log into your app across multiple platforms \cite{Facebook}.

\medskip

\noindent Option 3 - PHP \& SQL \\
\noindent Using PHP and SQL to compare submitted usernames and passwords against stored data on a database.

\medskip

\subsubsection{Goals}
\noindent An efficient and secure method for allowing for users to login and continue editing their documents from any computer or device they choose.


\medskip

\subsubsection{Evaluation Criteria}
\noindent The options are evaluated on
\begin{itemize}
\item Ease of Use
\item Features
\item Security
\end{itemize}


\newpage
\subsubsection{Option Comparison} ~\\

\tablehead{}
\begin{supertabular}{|p{2.5cm}|p{4.5cm}|p{4.5cm}|p{4.5cm}|}
\hline
	& OpenID
	& Facebook Login
	& PHP \& SQL\\
\hline
	Ease of Use
	& Requires Credentials with Corresponding Login Providers, Lots of available libraries
	& Requires Credentials with Facebook and an App ID with Facebook
	& Easy to implement, but if setup incorrectly can lead to problems\\
\hline
	Security
	& Relies on credential host and user security
	& Relies on Facebook and user security
	& Relies on password protection implementation and user security\\
\hline
	Features
	& Easily log in with OpenID partner credential hosts (Google, Microsoft, Yahoo, etc)
	& Easily log in with Facebook credentials
	& Easily log in with user created account and password \\
\hline

\end{supertabular}

\medskip

\subsubsection{Discussion}
\noindent The ideal user authentication system would be a combination of all three of the above implementations. While logging in with Facebook would make it easier to determine who is using the service, preventing unauthorized users from accessing unreleased copyrighted material, not everyone has a Facebook. Additionally using an OpenID login system would be reliant on other platform holders that use OAuth 2.0. Finally, using a self created account is often the easiest and can allow users to not be tied to a given account and also prevent private information from being retrieved from user accounts.

\medskip
\subsubsection{Selection}
\noindent For our implementation, we plan on using initially a PHP and SQL system to validate user account information on our database. Additionally, we will look into adding both OpenID and Facebook Login down the road.







\newpage
\subsection{Interface Design}
\noindent Option 1 - User Centered Design \\
\noindent A deep understanding of the target audience is able to provide insights into how to design and develop your application to suit your intended users \cite{Usability}.

\medskip

\noindent Option 2 - Activity-Centered Design  \\
\noindent Instead of focusing on research about intended users, the design and development are focused around making a given activity logically designed \cite{AListApart}.

\medskip

\noindent Option 3 - Self Design \\
\noindent The designer is responsible for representing the target audience. Though this can be a poor representation of the intended audience \cite{AListApart}.

\medskip

\subsubsection{Goals}
\noindent A design principle that allows for user interfaces that lead to user interfaces that are accepted by users and are easy to understand.

\medskip

\subsubsection{Evaluation Criteria}
\noindent The options are evaluated on
\begin{itemize}
\item Ease of Use
\item Strengths
\item Weaknesses
\end{itemize}

\newpage
\subsubsection{Option Comparison} ~\\

\tablehead{}
\begin{supertabular}{|p{2.5cm}|p{4.5cm}|p{4.5cm}|p{4.5cm}|}
\hline
	& User Centered Design
	& Activity-Centered Design
	& Self Design\\
\hline
	Ease of Use
	& Long process, that takes a lot of data gathering to provide insights into a target audience.
	& Easier to design an activity when not trying to cater a specific audience.
	& Very easy to design what works well for you as a designer. \\
\hline
	Strengths
	& Allows for the designer to understand what makes a user think the way they do.
		This allows for an interface design to be molded to an expected user.
	& Allows the designer to design a user interface based around an activity that
		a user will be performing instead of designing to a users wants and desires.
	& Allows for easy creation of user interface of how the designers see fit.
		Perfect for a target audience that is just like the designer. \\
\hline
	Weaknesses
	& Takes a long time to gather enough information to be able to design a
		good solution that feels natural to a target audience user.
	& Designed interface might work well for an intended activity, but could
		be antagonistic to a target audience.
	& Interface could be intended for an entirely different audience,
		leaving a confusing experience.\\
\hline

\end{supertabular}

\medskip

\subsubsection{Discussion}
\noindent The MVP Platform is highly user focused, which initially led the team to decided on User Centered Design. Activity-Centered Design or Self Design would greatly reduce the burden of research and discovery into what our target audience would like to see or be comfortable with naturally. Activity-Centered Design, if performed properly would result in interfaces that clearly work as intended, though might be off putting to our users. Self Design would allow for one of the team members to decide how a certain element shall look, but again can fall into an interface that does not satisfy our users.

\medskip
\subsubsection{Selection}
\noindent For our implementation, we plan on using User Centered Design. This is because users are our very important for our project. If our users do not like our user interface, then they will be less likely to use our software.








%
\newpage
\section{Josh Matteson}


\subsection{Testing}
\noindent Option 1 - Mocha \\
\noindent Mocha is a JavaScript testing framework, loaded with features. It runs on Node.js and also in the browser, making asynchronous testing simple and easy to use. Mocha tests run serially, allowing for flexible and accurate reporting, while mapping uncaught exceptions to the correct test cases. \cite{Mocha} \\

\medskip

\noindent Option 2 - QUnit \\
\noindent QUnit is a powerful, easy-to-use JavaScript unit testing framework. It's used by the jQuery, jQuery UI and jQuery Mobile projects and is capable of testing any generic JavaScript code. \cite{QUnit}  \\

\medskip

\noindent Option 3 - Jasmine \\
\noindent Jasmine is a behavior-driven development framework for testing JavaScript code. It does not depend on any other JavaScript frameworks. It does not require a DOM. And it has a clean, obvious syntax so that you can easily write tests. \cite{Jasmine}  \\

\subsubsection{Goals}
\noindent Using this technology will aid in proper functionality and minimize errors. Without properly testing code, a number of problems can occur that can disrupt and slow down progress in a team. In extreme cases, not properly testing could lead to failure of the application.\\

\subsubsection{Evaluation Criteria}
\noindent The options are evaluated on
\begin{itemize}
\item Ease of Use
\item Features
\item Documentation
\item Integration
\item Other
\end{itemize}


%\begin{center}

\newpage
\subsubsection{Option Comparison} ~\\
\tablehead{}
\begin{supertabular}{|p{2.5cm}|p{4.5cm}|p{4.5cm}|p{4.5cm}|}
\hline
	& Mocha
	& QUnit
	& Jasmine\\
\hline
	Ease of Use
	& Language is like spoken language, which makes for easier to understand
	& Not as friendly, language is similiar to other basic Unit testing frameworks
	& Language is like spoken English, fairly easy to understand\\
\hline
	Features
	& Spys, mocks, stubs, callbacks, etc. Most features
	& Most testing is done through assertions only, not as many features
	& Spys, mocks, stubs, callbacks, etc. Lacks Assertions, but can be implemented with library\\
\hline
	Documentation
	& Fairly common, moderate amount of documentation
	& Not well documented
	& Most popular, abundant documentation \\
\hline
	Integration
	& Moderate difficulty to set up
	& Easiest to set up
	& Moderate difficulty to set up\\
\hline
	Other
	&
	&
	& Familiar\\
\hline
\end{supertabular}

\medskip

\subsubsection{Discussion}
\noindent It is well known that an application that ensures proper functionality supersedes other applications and,
more importantly, competitors. Quality assurance is acknowledged as highly distinguished, and therefore,
an attribute deserving of notable consideration. Having this in mind, we will be considering these mentioned
testing frameworks: Mocha, QUnit, and Jasmine.


\subsubsection{Selection}
\noindent Jasmine, even though more difficult to integrate and set up, would be most feasible. Jasmine has many features that win
out over Mocha and QUnit, as well as thorough documenation. The test are easily grouped through describe blocks, which makes it easy
to identify where certain problems are. There is familiarity with Jasmine over the other testing frameworks, which means less time learning
for the team.

%\end{center}

\medskip







\newpage
\subsection{Revision Control Software}

{\noindent Option 1 - Git \par}
{\noindent Git is an open source, free, distributed version control system. It's designed to handle small to very large projects with speed and accuracy.
It is commonly used by a variety of companies. \cite{Git} \par}

\medskip
{\noindent Option 2 - CVS - Concurrent Versions System \par}
{\noindent CVS is the most popular and widely adopted revision control system to date. It is widely considered because of its low learning curve. \cite{JaxEnter}  \par}

\medskip
{\noindent Option 3 - SVN \par}
{\noindent Subversion is another open source version control software used by platforms like Ruby, Python Apache, and more. SVN has many versions and IDEs available. \cite{SVN} \par}

\medskip
\subsubsection{Goals}
{\noindent Our application will widely revolve around the revision/version control software that we use, which makess choicing the correct one ideal.
The main goals for our revision control would be ease of incorporation into our app, ease of use, and the features it comes with. More criteria found on the next page. \par}


\subsubsection{Evaluation Criteria}
\noindent The options are evaluated on
\begin{itemize}
\item Established
\item Features
\item Ease of Use
\item Speed
\end{itemize}


\subsubsection{Option Comparison} ~\\
\tablehead{}
\begin{supertabular}{|p{2.5cm}|p{4.5cm}|p{4.5cm}|p{4.5cm}|}
\hline
	\cite{JaxEnter}
	& Git
	& CVS
	& SVN\\
\hline
	Established
	& More than a decade
	& More than three decades
	& Around 15 years\\
\hline
	Features
	& Access to tree offline, highly distributed peer-to-peer model
	& Basic Features
	& Atomic Operations, Various IDE Plugins \\
\hline
	Ease of Use
	& Steep learning curve
	& Relatively low learning curve
	& Relatively low learning curve \\
\hline
	Speed
	& Very fast
	& Very slow
	& Slow \\
\hline
\end{supertabular}

\medskip

\subsubsection{Discussion}
{\noindent Revision control will be used frequently in our app to save progress and ensure progress hasn't been lost.
This in mind, we need software that ensures speed and can perform various operations inexpensively.
While most of the revision control software offers features, only one of them offers the speed required for frequent revisions. \par}

\medskip
\subsubsection{Selection}
{\noindent Although accompanied by a steep learning curve, Git offers all the features and speed required for the operations of our application.
Fortunately, most of our team members are familiar with using Git to one extent or another. This makes the trade off feasible. \par}







\newpage
\subsection{Database}
\noindent Option 1 - SQL Server \\
\noindent SQL Server is owned by Microsoft, and is known for its scalable, hybrid database platform. Being owned by Microsoft, this database is well known. \cite{SQLServer}

\medskip
\noindent Option 2 - MongoDB \\
\noindent MongoDB is a document based database, with an Expressive Query Language and Secondary Indexes. \cite{MongoDB}

\medskip
\noindent Option 3 - MySQL \\
\noindent Many of the worlds largest companies rely on MySQL, such as Facebook, Google, and others. It is a relational database. \cite{MySQL}

\subsubsection{Goals}
\noindent Having an understandable, well built database can aid in the flow of building an application as well as delivery. This in mind, we want to balance cost with effectiveness.


\subsubsection{Evaluation Criteria}
\noindent The options are evaluated on
\begin{itemize}
\item Cost
\item Type
\item Ease of Use
\end{itemize}



\subsubsection{Option Comparison} ~\\
\tablehead{}
\begin{supertabular}{|p{2.5cm}|p{4.5cm}|p{4.5cm}|p{4.5cm}|}
\hline & SQL Server & MongoDB & MySQL\\ \hline
Cost & Enterprise: \$12,256, Standard: \$3717, Developer: Free & OpenSource, but essential periphery software can range from \$45 a month to as much as \$5225 a month & Enterprise: \$5000, Standard: \$2000, Community Edition: Free \\ \hline
Type & Structured Query Language & Expressive Query Language & Structured Query Language\\ \hline
Ease of Use & Uses Queries to gather information & EQL data access and manipulation in sophisticated ways, operational and analytical applications \cite{ibmbpnetwork} & Uses Queries to gather information \\ \hline
\end{supertabular}

\medskip

\subsubsection{Discussion}
\noindent The most obvious and transparent qualities for a database are the cost and type, however, databases that come with a free edition allow for us to test and experiment while developing. This is imperative to the development process, and shouldn't be overlooked. Allowing for an easy transition from the free edition to a standard or express edition is important as well. A non relational database like MongoDB could be beneficial when not knowing the type of data that will be stored.


\subsubsection{Selection}
\noindent All things considered, SQL Server would be considered our top pick. Experience using an SQL based database before aids in the difficulty of using a query. Having a free edition contributes a moderate amount into the decision as well.








%
\newpage
\section{Evan Tschuy}

\subsection{Server Back-end}
\noindent Option 1 - NodeJS \\
\noindent NodeJS is a modern web back-end framework developed by the Node Foundation,
primarily led by Joyent. By using JavaScript its language of choice, Node allows
developers to use the language's unique concurrency paradigms to quickly develop
scalable applications.

\medskip

\noindent Option 2 - Django \\
\noindent The Django framework is a massive web framework developed in Python that comes "batteries
included". The Framework includes everything from geo-libraries to support for four different kinds of
databases, meaning a large initial learning curve but a large payoff.

\medskip

\noindent Option 3 - Flask \\
\noindent Flask is a micro-framework. It comes with the bare minimum needed to do HTTP handling, leaving
what other frameworks come with to an array of choices from third party developers. This means the core
framework is quick to learn, but can quickly leave a developer feeling constrained.

\medskip

\noindent Option 4 - Ruby on Rails \\
\noindent Ruby on Rails is the old standard of web frameworks. It was the original batteries included
framework, and has over the years been known for its ease of use. However, the framework is quite old
and shows some signs of age, using sometimes outdated paradigms and generally being less friendly to beginners than more modern frameworks.

\medskip

\subsubsection{Evaluation Criteria}
{\noindent The options are evaluated on

\begin{itemize}
\item Ease of use
\item Language
\item Features
\item Ecosystem
\item License
\end{itemize}

 \par}

%\begin{center}

\newpage
\subsubsection{Option Comparison} ~\\

\tablehead{}
\begin{supertabular}{|p{2.5cm}|p{3.3cm}|p{3.3cm}|p{3.3cm}|p{3.3cm}|}
\hline
	& NodeJS
	& Django
	& Flask
	& Ruby on Rails \\
\hline
	Ease of Use
	& Javascript backend shares language with frontend; super quick iteration
	& Large framework containing all needs within, with a large learning curve
	& Microframework containing only minimal needs; requires finding external packages
	& Old framework with massive, but aging, ecosystem \\
\hline
	Language
	& Javascript
	& Python
	& Python
	& Ruby \\
\hline
	Features
	& Minimal HTTP interaction, massive external ecosystem for extras
	& Massive built-ins
	& Minimal with external ecosystem for extras
	& Massive built-ins \\
\hline
	Ecosystem
	& Massively popular today with expansive and growing ecosystem
	& Decently large ecosystem with built-ins for most tasks
	& Decently large ecosystem
	& Large ecosystem but fairly old \\
\hline
	Release
	& 7.1.0
	& 1.10.3
	& 0.11
	& 5.0.0.1 \\
\hline
	License
	& MIT
	& BSD
	& BSD
	& MIT \\
\hline
\end{supertabular}

\medskip

\subsubsection{Discussion}
{\noindent
All backends listed are popular within their respective communities. However, more new projects are being created using NodeJS, as its modern paradigms and sharing of a language with frontend development allow developers familiar with it to iterate quicker and write more expressive code. Django's built-ins allow for a quicker initial development time but mean being isolated from the rapidly expanding ecosystem around NodeJS.
Ruby on Rails is a rather old backend, and has not shown to have the modern flexibility of Node.
 \par}

\medskip
\subsubsection{Selection}
{\noindent As NodeJS has the most expansive ecosystem, and allows us to share a common language between the front and backennds of the project, we will be using it over the other options considered. Additionally, Ward Cunningham's Federated Wiki uses it as one of its backends, and if we fork it, we can continue to use its NodeJS backend. \par}







\newpage
\subsection{Text Formatting}

\noindent Option 1 - Markdown \\
\noindent Markdown is a highly lightweight markup language that allows easy, human-readable markup
of text to include headings, bold/italic/underline/etc, bullets, and numbered lists.
 The original markdown does not support things like images or videos; Markdown has various
 "flavors", or implementations, that sometimes allow for such things.

\medskip

\noindent Option 2 - Restructured Text \\
\noindent Restructured Text is a markup language written in Python for writing documentation,
simple websites, etc. It allows for highly varied but still restricted markup; it allows
for image embeds, fancy linking, titles, etc. It does not allow users to embed arbitrary elements.

\medskip

\noindent Option 3 - Raw Latex \\
\noindent Storing simply raw Latex allows the greatest flexibility, as it is literally the
text that will be rendered by the backend. Raw Latex allows for things like image embeds,
math text, etc. but also raw include, and as such should be filtered to a
subset to be suitable for use in a public-facing scenario.

\medskip
\subsubsection{Evaluation Criteria}
{\noindent The options are evaluated on

\begin{itemize}
\item Ease of use for end-users
\item Markup options
\item Compile language
\item Security
\end{itemize}

 \par}

%\begin{center}

\newpage
\subsubsection{Option Comparison} ~\\

\tablehead{}
\begin{supertabular}{|p{2.5cm}|p{4.5cm}|p{4.5cm}|p{4.5cm}|}
\hline
	& Markdown
	& ReStructured Text
	& Latex \\
\hline
	Ease of Use
	& Easy to use, with minimal options and human-readable markup;
		different implementations have slight differences leading to confusion
	& Relatively human-readable markup but with massive number of options
	& Essentially infinite options but not very human-readable/human-writeable. Some users may already be proficient.  \\
\hline
	Markup
	& Options readable in single-page document, not allowing for high flexibility
	& Highly featureful with well-defined language
	& Allows for infinite options; controlling it might be difficult \\
\hline
	Compile Language
	& Dozens of different libraries for different languages,
		each with slightly different interpretation of the markup
	& Python
	& Compilation is needed to display but it is also the compilation being used for the final book, so the render is one-to-one \\
\hline
	Security
	& Small language leads to minimal exploits
	& Well-defined with real-world-tested libraries
	& Needs careful processing to stop end-users from inputting
		harmful raw input (being able to retrieve files from the host system)\\
\hline
\end{supertabular}

\medskip

\subsubsection{Discussion}
\noindent
All languages listed allow users to do simple things like bold text and link to other pages. However,
Latex offers the most flexibility and allows users to be able to do anything they want.
Allowing this while maintaining ease-of-use would suggest needing a way to prevent
users from including arbitrary paths on disk to scrape files such as private tokens.

\medskip

\subsubsection{Selection}
\noindent
The Federated Wiki uses HTML with minimal security processing. This shows that
insecure languages can still be used, with some precautions. As our books already
are rendered in Latex, it will be most reasonable to simply render text passed
in from users as Latex.










\newpage
\subsection{Password Storage}
\noindent Option 1 - Bcrypt \\
\noindent Bcrypt is a password hashing function that takes a very large amount of time to crack
an individual password -- it is designed to be slow. This means a hacker cannot simply crack
a database worth of passwords in one sitting, as with older hashes like MD5.

\medskip

\noindent Option 2 - Scrypt \\
\noindent Scrypt is designed to take up large amounts of time, and large amounts of RAM,
when hashing. This ensures that a hacker cannot simply buy a powerful CPU and crack passwords with
pure power. However, scrypt, being designed more-so for computer hard disk passwords, can take
multiple seconds and hundreds of megabytes of RAM to process.

\medskip

\noindent Option 3 - pbkdf2 \\
\noindent PBKDF2 is a function that repeatedly hashes a password using the HMAC, or "keyed-hash
message authentication code", function. For a CPU, cracking a large number of passwords using pbkdf2
is difficult, as it takes a large amount of time to crack an individual password. Using a GPU,
however, a large number of hashes can be run in parallel, making it quick to crack with high end hardware.

\medskip

\noindent Option 4 - raw storage \\
\noindent Another option for password storage is to store the passwords in plain text. This allows
users to recover their passwords directly through a password reminder email. However, this comes with
the major downside that compromising the database allows a hacker to be able to access any
accounts on unrelated services where users use the same username and password (a common pattern
in non-technical and technical users alike).

\medskip

\subsubsection{Evaluation Criteria}
{\noindent The options are evaluated on

\begin{itemize}
\item Cracking
\item Storage
\item Hacking Resistance
\end{itemize}

 \par}

\newpage
\subsubsection{Option Comparison} ~\\

\tablehead{}
\begin{supertabular}{|p{2.5cm}|p{3.3cm}|p{3.3cm}|p{3.3cm}|p{3.3cm}|}
\hline
	& Bcrypt
	& Scrypt
	& pbkdf2
	& plain text \\
\hline
	Cracking
	& Bcrypt is highly resistant to cracking on CPUs and GPUs,
  		but can be cracked quickly using specialized FPGAs.
	& Scrypt is highly resistant to cracking on CPUs, GPUs, etc but
  		takes a large amount of time to verify a valid password.
	& pbkdf2 is highly resistant to cracking on CPUs but can be easily cracked on a GPU.
	& A plaintext password does not need to be cracked as it is already stored as a raw password. \\
\hline
	Storage
	& can be stored in a database without issue.
	& can be stored in a database without issue.
	& can be stored in a database without issue.
  	& Plain passwords must be stored in a way that ensures they can
  		never be hacked, which is impossible. \\
\hline
	Hacking \hspace{3em} Resistance
  	& A bcrypt password hash must be cracked before it can be used elsewhere.
  	& Scrypt hashes must be cracked before use elsewhere.
  	& pbkdf2 hashes must be cracked before use elsewhere.
  	& A plain password, once retrieved from a database, can be used along with
  		the username/email associated to hack other sites. \\
\hline
\end{supertabular}

\medskip

\subsubsection{Discussion}
\noindent
Password storage is a tradeoff between ease of use and difficulty of reversing.
Scrypt is too slow for use on a website with many users, whereas plain text passwords are
too insecure as a hacker can reuse the password immediately without cracking. pbkdf2 and
bcrypt do a good job defending against CPU cracking, but as pbkdf2 can be cracked using a
GPU, bcrypt is left remaining as the best tradeoff between speed and cracking.


\medskip
\subsubsection{Selection}
\noindent
As mentioned above, bcrypt-hashed passwords present a good tradeoff between
cracking ability and verification speed. As such, the Many Voices Platform will use
bcrypt to securely verify any password used with the system.


\section{Conclusion}
\noindent  The Many Voices Publishing Platform is a combination of User Interfaces, Documentation,
User Centered Design, Testing, User Authentication, Databases, Server Back-end, Text Formatting,
Password Storage, and the users themselves. Determining the technologies behind these parts and
pieces is a difficult task to accomplish, as many choices can satisfy the requirements of the project.
Finding the best solution however is the goal of this document, to provide a clear path forward
for the platform as a whole.

\newpage
\bibliographystyle{IEEEtran}
\bibliography{techreview}


\end{document}
